\documentclass[10pt,letterpaper]{article}
\usepackage[margin=0.64in]{geometry}
\usepackage[T1]{fontenc}
\usepackage{lmodern,amsmath,amssymb,microtype}
\usepackage[colorlinks=true,urlcolor=blue]{hyperref}
\setlength{\parindent}{0pt}
\setlength{\parskip}{4pt}
\setlength{\abovedisplayskip}{5pt}
\setlength{\belowdisplayskip}{5pt}
\pagestyle{empty}
\newcommand{\lead}[1]{\par\smallskip\textbf{#1}\par\nobreak}
\begin{document}
\begin{center}
{\Large\bfseries Spin and charge response: the Hubbard TBC definition}\\[3pt]
{\small Definition and comparison with the student's response note\hfill 13 October 2026}
\end{center}
\vspace{-6pt}

\lead{1. Target observable at a fixed twist}
Let $H_\theta$ be an $L$-site Hubbard ring with a common spin-independent twist
$\theta$. Set the lattice spacing and $\hbar$ to one. The filling convention is
$x=N_e/(2L)$, so $x=0.5$ is half-filling. For a definite momentum transfer
$q_m=2\pi m/L$, define
\begin{equation}
 \rho_{c/s}(q_m)=\frac{1}{\sqrt L}\sum_{j=0}^{L-1}e^{iq_mj}
                  \bigl(n_{j\uparrow}\pm n_{j\downarrow}\bigr).
 \label{eq:rho}
\end{equation}
The plus/minus sign denotes charge/longitudinal spin. This spin convention uses
$2S^z$; using $S^z$ divides spin intensity by four, not excitation energies.
The zero-temperature \emph{inelastic} structure factor is
\begin{equation}
 S_\nu(q,\omega;\theta)=\sum_{E_a>E_0}
 |\langle a|\rho_\nu(q)|0\rangle|^2
 \delta\!\left[\omega-(E_a-E_0)\right],\qquad \nu=c,s.
 \label{eq:lehmann}
\end{equation}
Separate all zero-energy (elastic) weight before broadening. Equivalently, with
$\widetilde\rho_\nu=P_{\rm exc}\rho_\nu$ and $P_{\rm exc}$ excluding the entire
ground-state subspace,
\begin{equation}
 F_\nu(q,z)=\langle0|\widetilde\rho_\nu^\dagger(q)
 (z+E_0-H_\theta)^{-1}\widetilde\rho_\nu(q)|0\rangle,
 \qquad S_\nu^\eta=-\pi^{-1}\operatorname{Im}F_\nu(q,\omega+i\eta).
\end{equation}
$F$ is the positive-frequency resolvent, not the full retarded commutator
susceptibility, which also has a negative-frequency term. Both total charge and
longitudinal spin are conserved: $S_c(0,\omega>0)=S_s(0,\omega>0)=0$.

\lead{2. Distinction from the written note}
Equation (1) of the student's note has the appropriate Lehmann form once its
operator is specified. Equation (2), as written, is a complex resolvent;
obtaining a real spectral function requires $-\operatorname{Im}/\pi$.
For the physical density, contributions of compatible momentum pairs must be
summed \emph{coherently} before squaring:
\begin{equation}
 \rho_{c/s}(q_m)=\frac{1}{\sqrt L}\sum_{p,\sigma}s_\sigma^{c/s}
 c^\dagger_{p+q_m,\sigma}c_{p,\sigma},\qquad
 s_\sigma^c=1,\quad s_\uparrow^s=1,\ s_\downarrow^s=-1.
\end{equation}
Here $p$ runs over the ring's allowed twisted momenta. An orbital-pair sum
\emph{outside} the modulus, as printed in Eq.~(2), generally omits interference
terms and is not the same total-density observable. The common twist shifts
all $p$ equally and cancels from the transfer $q_m$.

\lead{3. Distinction from the supplied implementation}
The supplied code \texttt{spin\_charge\_v5.py} retains cross terms, but uses the
different vertex $e^{iq(\alpha-\alpha')}$ with
$c^\dagger_{\alpha'\sigma}c_{\alpha\sigma}$. In its orbital-site basis this gives
\begin{equation}
 R_{c/s}(q)=L\bigl(d^\dagger_{q\uparrow}d_{q\uparrow}
             \pm d^\dagger_{q\downarrow}d_{q\downarrow}\bigr),\qquad
 d_{q\sigma}=L^{-1/2}\sum_\alpha e^{iq\alpha}c_{\alpha\sigma}.
\end{equation}
This is a valid \emph{Fourier-mode occupation} response: $q$ labels a mode,
not a fixed transferred momentum. It is not a normalization or gauge variant
of Eq.~\eqref{eq:rho}. In the same site basis, the physical vertex instead is
$\Gamma_{\alpha'\alpha}(q)=\delta_{\alpha'\alpha}e^{iq\alpha}/\sqrt L$.
Restricting the existing phase-difference vertex to $\alpha=\alpha'$ alone would
give phase one and would not fix it. Finite-frequency $R(0,\omega)$ can be
nonzero; its peak slopes do not directly define physical spin/charge velocities.

\lead{4. Match the ensemble before comparing}
At each twist, match $L,t,U$, the particle/spin sector or chemical-potential
prescription, degenerate-ground-state weights, normalization, and $\eta$.
For neutral excitations within one sector, $\mu N_e$ cancels from $E_a-E_0$;
$\mu$ selects the sector. The supplied driver selects sectors using a common
$\mu$ and then sums over reduced-BZ $k$: this is \emph{not} one fixed-twist ring.
In its boundary-phase convention $\theta_k=Lk$; a normalized comparison uses
$\overline S_\nu=\sum_k w_k S_\nu(q,\omega;\theta_k)$ with $\sum_k w_k=1$.
Benchmark the physical operators block by block first, then apply identical
$k$/twist averaging. No solver rerun is needed to establish these definitions.

\vfill
{\footnotesize\textbf{Sources.} Standard site-density conventions:
H.~Benthien and E.~Jeckelmann,
\href{https://arxiv.org/abs/cond-mat/0606748}{cond-mat/0606748}, Eqs.~(3),(5).
Comparison: supplied response note, Eqs.~(1),(2), and supplied code,
\texttt{spin\_charge\_v5.py}, lines 64, 74--105, 171--173.
These are distinctions between observables, not a claim that the student's
Lehmann solver is numerically incorrect.}
\end{document}
